\chapter{Типизированное вложение хопфовского цикла в 43-мерную ячейку}
\label{ch:v10-hopf-cycle-k43-typed-embedding}

\section{Третья вершина цикла}

В $\mathcal K_{43}$ уже канонически различены вакуумная линия $|0\rangle$ и
гиперзарядовая линия $|Y\rangle$. Для трёхвершинного сквозного цикла нужна
третья линия. В старом шумовом репере имеются тридцать транспортных
направлений и двенадцать калибровочных направлений. Поэтому возникают две
естественные конструкции.

Первая выбирает одно транспортное направление $|T_a\rangle$ и задаёт
изометрию
\begin{equation}
 E_a=(|0\rangle,|Y\rangle,|T_a\rangle):
 \mathbb C^3\hookrightarrow\mathcal K_{43}.
 \label{eq:v10-hopf-coordinate-embedding}
\end{equation}
Она сохраняет генератор цикла:
\begin{equation}
 E_a^*L_{43}^{(a)}E_a=L_\circlearrowright,
 \qquad
 L_{43}^{(a)}=E_aL_\circlearrowright E_a^*.
 \label{eq:v10-hopf-coordinate-compression}
\end{equation}
Полный оператор имеет ранг $2$, ядро размерности $41$, сохраняет сумму
вероятностей и не содержит отрицательных внедиагональных скоростей.
Однако существует тридцать равноправных выборов $a$; прежняя архитектура не
выделяет один из них.

\section{Симметричная транспортная линия}

Вторая конструкция устраняет произвольный выбор посредством линии
\begin{equation}
 |T_{\rm sym}\rangle
 =\frac1{\sqrt{30}}\sum_{a=1}^{30}|T_a\rangle.
 \label{eq:v10-hopf-symmetric-transfer-line}
\end{equation}
Она инвариантна при перестановках транспортных меток, а вложение
$E_{\rm sym}=(|0\rangle,|Y\rangle,|T_{\rm sym}\rangle)$ остаётся
изометрическим. Сжатие снова точно воспроизводит абстрактный цикл.

Но обратное расширение на весь координатный вероятностный симплекс создаёт
между различными транспортными метками элементы
\begin{equation}
 (L_{43}^{\rm sym})_{ab}=-\frac{\kappa}{10}<0,
 \qquad a\ne b.
 \label{eq:v10-hopf-symmetric-negative-rate}
\end{equation}
Таких упорядоченных внедиагональных элементов ровно
$30\cdot29=870$. Следовательно, когерентная симметризация является хорошим
гильбертовым вложением, но не классическим марковским генератором на 43
координатных состояниях.

\section{Точная дихотомия}

Два требования --- марковская положительность и каноническая симметрия
третьей линии --- дают матрицу
\begin{equation}
 \begin{pmatrix}1&0\\0&1\end{pmatrix},
 \label{eq:v10-hopf-embedding-dichotomy}
\end{equation}
где первая строка относится к координатному вложению, вторая --- к
симметричному. Вектор полных проходов равен $(0,0)^T$.

\begin{theorem}[Алгебраическое вложение существует, каноническое марковское --- нет]
Любая отдельная транспортная метка даёт точное марковское вложение цикла в
$K_{43}$, но выбор имеет кратность $30$. Симметричная транспортная линия
устраняет кратность, однако создаёт $870$ отрицательных внедиагональных
скоростей. Ни одна из двух конструкций не удовлетворяет обоим требованиям.
\end{theorem}

\begin{nogocriterion}[Запрет когерентной линии как классической вероятности]
Изометрическое сжатие генератора не гарантирует положительность его
расширения на координатный вероятностный симплекс. Когерентную суперпозицию
транспортных меток нельзя объявлять третьим классическим состоянием без
отдельного квантового марковского описания или дополнительного типового
множителя.
\end{nogocriterion}

\begin{protocol}[Статус вложения цикла]\begin{enumerate}\raggedright
\item координатное марковское вложение: \textbf{получено};
\item число равноправных координатных выборов: \textbf{$30$};
\item симметричное изометрическое вложение: \textbf{получено};
\item отрицательных внедиагональных скоростей: \textbf{$870$};
\item полных проходов двух конструкций: \textbf{$0/2$};
\item архитектура аудита: \textbf{$8/8$};
\item каноническое марковское происхождение: \textbf{$0/1$};
\item абсолютный масштаб скорости: \textbf{$0/1$};
\item следующий гейт: \textbf{вложение цикла в произведение $K_{43}$ с КМС-носителем}.
\end{enumerate}\end{protocol}

Машинный сертификат сохранён в
\path{s2t_v10_cell_birth_four_volume_hopf_cycle_k43_typed_embedding_gate_results.json}.